% This file is part of babel. For further details see:
% https://www.ctan.org/pkg/babel
\ifx\BabelBeforeIni\undefined
  \PackageError{babel}%
    {This file is a component of babel and cannot\MessageBreak
     be loaded directly. I'll stop immediately}%
    {Just use babel as documented.}%
  \stop
\fi
\BabelBeforeIni{az}{%
  \ifx\Azerbaijanischwa\undefined
    % Azerbaijani Latin schwa (Ə U+018F, ə U+0259). T1 has no glyph for
    % this Latin letter; the default definition uses the raw codepoint
    % via \char (not the literal character) so it is safe even where
    % \DeclareUnicodeCharacter later maps that same codepoint back to
    % this very macro -- using the literal character there would make
    % that mapping self-referential and loop forever. \char with a
    % Unicode codepoint is resolved directly by LuaTeX/XeTeX's native
    % Unicode font access, so this remains correct on those engines.
    % Under pdfTeX/T1 the visual glyph is instead borrowed, via the
    % standard \UseTextSymbol mechanism, from whichever Cyrillic
    % encoding (T2A/T2B/T2C/X2 -- all four place it at the same slot:
    % 154 upper, 186 lower) the user has also loaded; the macro's own
    % Latin identity is unaffected by that borrowing. This follows the
    % historical babel-azerbaijani package (Namig J. Guliyev).
    \DeclareTextCommandDefault{\Azerbaijanischwa}{\char"18F\relax}%
    \DeclareTextCommandDefault{\azerbaijanischwa}{\char"259\relax}%
    % \DeclareUnicodeCharacter is pdfTeX-kernel-only (it is undefined
    % under LuaTeX/XeTeX, which read UTF-8 natively and need no such
    % registration); guard it so loading this file never errors there.
    \ifx\DeclareUnicodeCharacter\undefined\else
      \DeclareUnicodeCharacter{018F}{\Azerbaijanischwa}%
      \DeclareUnicodeCharacter{0259}{\azerbaijanischwa}%
    \fi
    % The T1/Cyrillic-borrow setup below is pdfTeX-only: LuaTeX and
    % XeTeX read UTF-8 natively and already render Ə/ə correctly (glyph
    % coverage permitting) via the \char-based default above, without
    % needing T1 or any Cyrillic encoding at all.
    \ifx\directlua\undefined\ifx\XeTeXrevision\undefined
      \AtBeginDocument{%
        \let\bbl@az@cyrenc\@empty
        \@ifundefined{T@T2C}{}{\def\bbl@az@cyrenc{T2C}}%
        \@ifundefined{T@T2B}{}{\def\bbl@az@cyrenc{T2B}}%
        \@ifundefined{T@T2A}{}{\def\bbl@az@cyrenc{T2A}}%
        \@ifundefined{T@X2}{}{\def\bbl@az@cyrenc{X2}}%
        \ifx\bbl@az@cyrenc\@empty
          \PackageError{babel}%
            {No Cyrillic encoding (T2A, T2B, T2C or X2) is loaded.\MessageBreak
             Azerbaijani Latin schwa (Ə/ə) cannot be represented\MessageBreak
             under pdfTeX using T1 alone. Load one together with T1,\MessageBreak
             e.g. \string\usepackage[T1,T2A]{fontenc}}%
            {See the Azerbaijani section of the babel manual.}%
        \else
          \DeclareTextCommand{\Azerbaijanischwa}{T1}{%
            \UseTextSymbol{\bbl@az@cyrenc}{\char154}}%
          \DeclareTextCommand{\azerbaijanischwa}{T1}{%
            \UseTextSymbol{\bbl@az@cyrenc}{\char186}}%
        \fi}%
    \fi\fi
  \fi
}
\endinput
